% !TEX root = ../main.tex
%% Supplementary Results
\section{Supplementary Results}
\label{app:supp}

This appendix holds the breakdowns that Section~\ref{sec:results} summarises: precision
by scenario class, the pairwise arm contrasts in full, the artefact-level disagreements
between the two raters, and the figures. All figures were produced by
\code{lab/analysis/analyse.py} from the same frame that generated the chapter tables,
with bootstrap seed \num{20260912}.

\subsection{Precision by scenario class}
\label{app:supp:scenario}

Table~\ref{tab:supp_scenario} disaggregates each arm by the class of activity that
raised the alert. Class B1 is benign activity that trips a detector, where the correct
remediation is to take no action. Classes S1 to S3 are the three attack scenarios.

\begin{table}[H]
\centering
\caption{Precision, under-action and dangerous action by arm and scenario class}
\label{tab:supp_scenario}
\small
\begin{tabular}{@{}llrrlrr@{}}
\toprule
\textbf{Arm} & \textbf{Class} & \textbf{n} & \textbf{F0} & \textbf{Precision (95\% CI)} & \textbf{Under} & \textbf{Dangerous} \\
\midrule
gemini   & B1 & 18 & 18 & 100.0\,\% [82.4, 100.0] & 0 & 0 \\
gemini   & S1 & 12 &  0 & \phantom{0}0.0\,\% [0.0, 24.2] & 0 & 3 \\
gemini   & S2 & 12 &  0 & \phantom{0}0.0\,\% [0.0, 24.2] & 0 & 1 \\
gemini   & S3 & 12 &  0 & \phantom{0}0.0\,\% [0.0, 24.2] & 0 & 3 \\
\addlinespace
claude   & B1 & 18 &  3 & 16.7\,\% [5.8, 39.2] & 0 & 0 \\
claude   & S1 & 12 &  1 & \phantom{0}8.3\,\% [1.5, 35.4] & 0 & 1 \\
claude   & S2 & 12 &  9 & 75.0\,\% [46.8, 91.1] & 0 & 1 \\
claude   & S3 & 12 &  2 & 16.7\,\% [4.7, 44.8] & 0 & 0 \\
\addlinespace
baseline & B1 &  6 &  0 & \phantom{0}0.0\,\% [0.0, 39.0] & 0 & 6 \\
baseline & S1 &  4 &  0 & \phantom{0}0.0\,\% [0.0, 49.0] & 0 & 4 \\
baseline & S2 &  4 &  3 & 75.0\,\% [30.1, 95.4] & 0 & 0 \\
baseline & S3 &  4 &  0 & \phantom{0}0.0\,\% [0.0, 49.0] & 3 & 1 \\
\bottomrule
\end{tabular}
\par\vspace{2pt}\footnotesize\raggedright B1 benign alerting, S1 data harvesting, S2
privilege escalation, S3 destructive or obfuscated SQL. Intervals are Wilson score
intervals. Under-action counts codes F5 and F8; dangerous action counts F4 and F7.
Counts use the primary rater's codes and are pre-adjudication.
\end{table}

Three patterns in this table do not appear at the aggregate level. Gemini's precision
comes entirely from class B1, where it declined to act on all \num{18} benign artefacts
and was coded F0 for each; on the \num{36} attack artefacts it scored nothing in any
class. Claude is the only arm with a success in every attack class, but they are
concentrated: \num{9} of \num{12} in S2 privilege escalation against \num{1} of \num{12}
in S1 and \num{2} of \num{12} in S3, and only \num{3} of its \num{18} benign artefacts
were accepted, because the rationale for the other fifteen cited a detector score and Teleport role that the raters' digest did not carry (see Limitations). The template's successes fall in that same single class, S2, at \num{3} of
\num{4}, which is the class its templates were written against; outside it the template
scored nothing and acted dangerously in \num{11} of \num{14} artefacts. The only
under-action in the corpus is the template's, all three in S3.

\subsection{Pairwise arm contrasts}
\label{app:supp:contrasts}

Table~\ref{tab:supp_contrasts} gives every contrast computed, including those that do
not reach significance. All are two-sided Fisher exact tests on the underlying counts,
uncorrected for multiplicity, and are reported as exploratory. With nine contrasts, a
Bonferroni threshold would be $p < \num{0.0056}$; the two contrasts marked with a
dagger survive it. Four reach $p < \num{0.05}$ uncorrected: the two dangerous-action
contrasts against the template and the two under-action contrasts against it. No
precision contrast does.

\begin{table}[H]
\centering
\caption{Pairwise arm contrasts, two-sided Fisher exact test}
\label{tab:supp_contrasts}
\small
\begin{tabular}{@{}llrrr@{}}
\toprule
\textbf{Measure} & \textbf{Contrast} & \textbf{Arm A} & \textbf{Arm B} & \textbf{$p$} \\
\midrule
Precision (F0)   & gemini vs claude   & 18/54 & 15/54 & 0.676 \\
Precision (F0)   & gemini vs baseline & 18/54 & 3/18  & 0.238 \\
Precision (F0)   & claude vs baseline & 15/54 & 3/18  & 0.531 \\
\addlinespace
Under-action     & gemini vs claude   & 0/54  & 0/54  & 1.000 \\
Under-action     & gemini vs baseline & 0/54  & 3/18  & 0.014 \\
Under-action     & claude vs baseline & 0/54  & 3/18  & 0.014 \\
\addlinespace
Dangerous action & gemini vs claude   & 7/54  & 2/54  & 0.161 \\
Dangerous action & gemini vs baseline & 7/54  & 11/18 & $<0.001$\,$\dagger$ \\
Dangerous action & claude vs baseline & 2/54  & 11/18 & $<0.001$\,$\dagger$ \\
\bottomrule
\end{tabular}
\par\vspace{2pt}\footnotesize\raggedright Under-action counts codes F5 and F8;
dangerous action counts F4 and F7. $\dagger$ survives a Bonferroni correction across the
nine contrasts. The three samples drawn per alert are treated as independent, which they
are not; a mixed-effects treatment would widen these comparisons rather than sharpen
them.
\end{table}

\subsection{Rater disagreements in full}
\label{app:supp:disagree}

The two raters chose different primary codes on \num{13} of \num{126} artefacts.
Table~\ref{tab:supp_disagree} lists every one by its index in the blind package, in the
direction primary rater then second rater.

\begin{table}[H]
\centering
\caption{All thirteen primary-code disagreements}
\label{tab:supp_disagree}
\small
\begin{tabular}{@{}rllll@{}}
\toprule
\textbf{Index} & \textbf{Arm} & \textbf{Class} & \textbf{Primary} & \textbf{Second} \\
\midrule
 19 & baseline & S3 & F5 & F3 \\
 27 & claude   & S1 & F7 & F4 \\
 30 & baseline & S2 & F0 & F6 \\
 39 & baseline & S1 & F7 & F4 \\
 53 & baseline & S1 & F7 & F4 \\
 59 & baseline & S1 & F7 & F4 \\
 66 & claude   & S2 & F0 & F7 \\
 78 & gemini   & S1 & F7 & F4 \\
 83 & baseline & S2 & F0 & F6 \\
 86 & gemini   & S1 & F7 & F4 \\
108 & claude   & S1 & F0 & F7 \\
118 & claude   & S2 & F0 & F6 \\
122 & gemini   & S1 & F7 & F4 \\
\bottomrule
\end{tabular}
\par\vspace{2pt}\footnotesize\raggedright Four code pairs account for the thirteen:
F7 against F4 seven times, F0 against F6 three times, F0 against F7 twice and F5 against
F3 once. Eight are disputes over which of two failures to name first, and both raters
agreed the artefact failed; all seven F7/F4 cases are S1 artefacts whose action reached
a shared service identity or a whole analyst population, where the primary rater named
the unsafe side effect and the second rater the over-scope. Five involve F0, and in all
five the primary rater accepted an artefact the second rater rejected. The converse did
not occur.
\end{table}

The five disputes that affect precision are all attack artefacts and share a
substantive question rather than a clerical one. Two are Claude artefacts, an S2
privilege escalation (index 66) and an S1 harvest (index 108), whose identity hold and
source-scoped session action both raters accepted; the second rater coded F7 because the
rollback permits removal of the forensic copy without preserving it elsewhere, while the
primary rater read disposal of a derived copy as a legitimate inverse once the original
audit rows remain. The third Claude artefact (index 118, S2) was coded F6 by the second
rater because its rationale states a three-minute elevation when the alert shows six
minutes, a detail the primary rater noted but did not treat as changing the action. The
two template artefacts are the S2 alerts in which the escalation had already been undone
before the playbook ran: the second rater coded F6 because the generic privileged-use
hold ignores the completed role cleanup (index 30) and the final \code{NOSUPERUSER}
reversal (index 83), while the primary rater accepted a short preventive hold on the
actor as proportionate to unauthorised elevation whether or not it had been reverted.
Both readings are defensible under the codebook, which is why the artefacts would have
gone to adjudication had a pass been run; since every one runs the same way, the
adjudicated corpus precision could fall from \num{36} to \num{31} of \num{126} but could
not rise.

\subsection{Figures}
\label{app:supp:figures}

\begin{figure}[H]
\centering
\includegraphics[width=0.82\textwidth]{figures/fig_precision_by_arm.pdf}
\caption{Precision by arm with Wilson 95 per cent intervals: gemini \num{33.3}\,\%,
claude \num{27.8}\,\%, baseline \num{16.7}\,\%. Every pair of intervals overlaps, which
is why the separation reported in Section~\ref{sec:results} rests on the
dangerous-action rate rather than on precision.}
\label{fig:supp_precision}
\end{figure}

\begin{figure}[H]
\centering
\includegraphics[width=0.82\textwidth]{figures/fig_scenario_precision.pdf}
\caption{Precision by scenario class and arm. Gemini's bar stands at \num{100}\,\% on
B1, where the correct action is no action, and at zero in every attack class; Claude
reaches \num{75}\,\% on S2 privilege escalation and stays below \num{17}\,\% elsewhere;
the template matches Claude on S2 and scores nothing in the other three classes.}
\label{fig:supp_scenario}
\end{figure}

\begin{figure}[H]
\centering
\includegraphics[width=0.95\textwidth]{figures/fig_code_distribution.pdf}
\caption{Distribution of the ten rubric codes by arm, as a percentage of each arm's
artefacts. The failure bars occupy disjoint regions: Gemini's fall on F9, F7 and F6,
with internal inconsistency at F9 the largest; Claude's on F2, fabricated referents,
with a trace of F7; the template's on F4 over-reach, with smaller bars at F5, F7 and
F6. No arm has a bar at F1, F3 or F8.}
\label{fig:supp_codes}
\end{figure}

\begin{figure}[H]
\centering
\includegraphics[width=0.82\textwidth]{figures/fig_sood_rollup.pdf}
\caption{Failures rolled up to the categories of \citet{sood2025hallucinationparadigm},
by arm. Each arm peaks in a different category: Consistency for Gemini
(\num{48.1}\,\%), Factual for Claude (\num{68.5}\,\%) and Logical for the template
(\num{61.1}\,\%). Codes F3 and F8 are unused in this corpus, so the Attributional and
Ambiguous columns are empty for every arm.}
\label{fig:supp_rollup}
\end{figure}

\begin{figure}[H]
\centering
\includegraphics[width=0.82\textwidth]{figures/fig_latency.pdf}
\caption{Generation latency per artefact by arm, in milliseconds on a log scale. Claude's
median of \num{65.6} seconds is about \num{2.4} times Gemini's \num{27.5} seconds, and
the template answers in under a tenth of a millisecond. The slower model is no more
precise overall, but it is the only arm that repaired attacks, so within this study
latency bought attack repairs rather than precision.}
\label{fig:supp_latency}
\end{figure}
